\begin{lemma}
\label{lemma-lift-continuous}
Let $R \to S$ be a ring map. Let $\mathfrak n$ be an ideal of $S$.
Assume that $R \to S$ is formally smooth in the $\mathfrak n$-adic
topology. Consider a solid commutative diagram
$$
\xymatrix{
S \ar[r]_\psi \ar@{-->}[rd] & A/J \\
R \ar[r] \ar[u] & A \ar[u]
}
$$
of homomorphisms of topological rings where $A$ is adic
and $A/J$ is the quotient (as topological ring) of $A$ by a closed ideal
$J \subset A$ such that $J^t$ is contained in an ideal of definition
of $A$ for some $t \geq 1$. Then there exists a dotted arrow in the category of
topological rings which makes the diagram commute.
\end{lemma}

\begin{proof}
Let $I \subset A$ be an ideal of definition so that $I \supset J^t$
for some $t$. Then $A = \lim A/I^n$ and $A/J = \lim A/J + I^n$
because $J$ is assumed closed. Consider the following diagram of
discrete $R$ algebras $A_{n, m} = A/J^n + I^m$:
$$
\xymatrix{
A/J^3 + I^3 \ar[r] \ar[d] &
A/J^2 + I^3 \ar[r] \ar[d] &
A/J + I^3 \ar[d] \\
A/J^3 + I^2 \ar[r] \ar[d] &
A/J^2 + I^2 \ar[r] \ar[d] &
A/J + I^2 \ar[d] \\
A/J^3 + I \ar[r] &
A/J^2 + I \ar[r] &
A/J + I
}
$$
Note that each of the commutative squares defines a surjection
$$
A_{n + 1, m + 1} \longrightarrow A_{n + 1, m} \times_{A_{n, m}} A_{n, m + 1}
$$
of $R$-algebras whose kernel has square zero.
We will inductively construct $R$-algebra maps
$\varphi_{n, m} : S \to A_{n, m}$.
Namely, we have the maps $\varphi_{1, m} = \psi \bmod J + I^m$.
Note that each of these maps is continuous as $\psi$ is.
We can inductively choose the maps $\varphi_{n, 1}$ by starting
with our choice of $\varphi_{1, 1}$ and lifting up, using the
formal smoothness of $S$ over $R$, along the bottom row of the
diagram above. We construct the remaining maps $\varphi_{n, m}$
by induction on $n + m$. Namely, we choose $\varphi_{n + 1, m + 1}$
by lifting the pair $(\varphi_{n + 1, m}, \varphi_{n, m + 1})$
along the displayed surjection above (again using the formal smoothness
of $S$ over $R$). In this way all of the maps $\varphi_{n, m}$ are
compatible with the transition maps of the system.
As $J^t \subset I$ we see that for example
$\varphi_n = \varphi_{nt, n} \bmod I^n$ induces a map $S \to A/I^n$.
Taking the limit $\varphi = \lim \varphi_n$ we obtain a map
$S \to A = \lim A/I^n$. The composition into $A/J$ agrees
with $\psi$ as we have seen that $A/J = \lim A/J + I^n$.
Finally we show that $\varphi$ is continuous. Namely, we know that
$\psi(\mathfrak n^r) \subset J + I/J$ for some $r \geq 1$ by our
assumption that $\psi$ is a morphism of topological rings, see
Lemma \ref{lemma-continuous}. Hence $\varphi(\mathfrak n^r) \subset J + I$
hence $\varphi(\mathfrak n^{rt}) \subset I$ as desired.
\end{proof}